% BHU PYQ -- Transcribed & Corrected
% Paper: BHUVA-13: Ayurveda -- Traditional Health Care & Management Principles (Sem I Mid Term)
% Exam: Under Graduate Semester I Mid-Term Examination, 2024-25 / 2025-26 (NEP)
% Category: Value Added Course (VAC)

\documentclass[11pt,a4paper]{article}
\usepackage[a4paper,top=2.5cm,bottom=2.5cm,left=2.8cm,right=2.8cm,headheight=14pt]{geometry}
\usepackage{amsmath,amssymb}
\usepackage[T1]{fontenc}
\usepackage[utf8]{inputenc}
\usepackage{lmodern,microtype,fancyhdr,enumitem,hyperref,lastpage,parskip}

\hypersetup{
    colorlinks=true,
    urlcolor=black,
    linkcolor=black,
    pdftitle={BHUVA13: Ayurveda Sem I Mid-Term -- BHU 2024-25 (NEP VAC)}
}

\pagestyle{fancy}
\fancyhf{}
\renewcommand{\headrulewidth}{0.4pt}
\renewcommand{\footrulewidth}{0.4pt}
\fancyhead[L]{\small\textit{Banaras Hindu University}}
\fancyhead[R]{\small\textit{BHUVA-13: Ayurveda (Sem I Mid Term VAC)}}
\fancyfoot[C]{\small Page \thepage\ of \pageref{LastPage}}

\newlist{parts}{enumerate}{2}
\setlist[parts,1]{label=(\Alph*),leftmargin=*,itemsep=3pt,topsep=2pt}
\newcommand{\pts}[1]{\hfill\textit{[#1]}}

\begin{document}

\begin{center}
  {\large\bfseries BANARAS HINDU UNIVERSITY}\\[2pt]
  {\normalsize Under Graduate --- Semester I Mid-Term Examination, 2024-25 / 2025-26 (NEP)}\\[6pt]
  \rule{0.8\linewidth}{0.5pt}\\[6pt]
  {\large\bfseries Value Added Course (VAC)}\\[4pt]
  {\large\bfseries Paper Code: BHUVA-13 : Ayurveda --- Traditional Health Care \& Management Principles (Sem I Mid Term)}\\[4pt]
  {\normalsize Time: 1 Hour\hfill Maximum Marks: 30}
\end{center}

\rule{\linewidth}{0.4pt}

\smallskip
\noindent\textit{Instructions: Answer all 30 Multiple Choice Questions. Each question carries equal marks (1 mark each).}

\bigskip

\noindent\textbf{Question 1.} Maharshi Patanjali is mainly related to:\pts{1}
\begin{parts}
  \item Yoga
  \item Ayurveda
  \item Nyaya
  \item Jyotish
\end{parts}

\medskip

\noindent\textbf{Question 2.} Maheshwar Sutras are primarily related to:\pts{1}
\begin{parts}
  \item Grammar / Sanskrit Philology
  \item Philosophy
  \item Ayurvedic treatment
  \item Astrology
\end{parts}

\medskip

\noindent\textbf{Question 3.} The classical text ``Sushruta Samhita'' was written by:\pts{1}
\begin{parts}
  \item Maharshi Sushruta
  \item Maharshi Atreya
  \item Maharshi Charaka
  \item Maharshi Vagbhata
\end{parts}

\medskip

\noindent\textbf{Question 4.} Synonyms of ``Ayu'' (Life) in Ayurveda include:\pts{1}
\begin{parts}
  \item Dhari, Jivita, Nityaga, and Anubandha
  \item Only Dhari
  \item Only Anubandha
  \item None of the above
\end{parts}

\medskip

\noindent\textbf{Question 5.} Which of the following are counted in the Shad Padarthas of Vaisheshika Philosophy?\pts{1}
\begin{parts}
  \item Dravya, Guna, Karma, Samanya, Vishesh, and Samavaya
  \item Samanya \& Vishesh only
  \item Karma \& Samavaya only
  \item All of the above
\end{parts}

\medskip

\noindent\textbf{Question 6.} Which of the following is counted as an Aastika (Orthodox) Darshana?\pts{1}
\begin{parts}
  \item Nyaya, Vaisheshika, Samkhya, Yoga, Mimamsa, and Vedanta
  \item Charvaka
  \item Bauddha
  \item Jaina
\end{parts}

\medskip

\noindent\textbf{Question 7.} Total number of Sharirik Doshas in Ayurveda are:\pts{1}
\begin{parts}
  \item 3 (Vata, Pitta, Kapha)
  \item 2
  \item 4
  \item 5
\end{parts}

\medskip

\noindent\textbf{Question 8.} Total number of Manasik Doshas are:\pts{1}
\begin{parts}
  \item 2 (Raja and Tama)
  \item 3
  \item 4
  \item 5
\end{parts}

\medskip

\noindent\textbf{Question 9.} Astangahrudaya was written by:\pts{1}
\begin{parts}
  \item Acharya Vagbhata
  \item Acharya Charaka
  \item Acharya Sushruta
  \item Acharya Bhavamishra
\end{parts}

\medskip

\noindent\textbf{Question 10.} Who are known as the Devabhishak (Physicians of the Gods)?\pts{1}
\begin{parts}
  \item Ashwini Kumaras
  \item Brahma
  \item Vishnu
  \item Shiva
\end{parts}

\medskip

\noindent\textbf{Question 11.} Swasthavritta in Ayurveda primarily promotes:\pts{1}
\begin{parts}
  \item Preventive health and positive lifestyle (Dinacharya \& Ritucharya)
  \item Surgical procedures
  \item Emergency treatment
  \item Diagnostic tests
\end{parts}

\medskip

\noindent\textbf{Question 12.} Rasa Shastra in Ayurveda mainly deals with:\pts{1}
\begin{parts}
  \item Processing of metals, minerals, and mercurial preparations
  \item Surgical procedures
  \item Treatment of newborns
  \item None of the above
\end{parts}

\medskip

\noindent\textbf{Question 13.} Chemical symbol of Parad (Mercury) in Rasa Shastra is:\pts{1}
\begin{parts}
  \item Hg
  \item Ag
  \item Au
  \item Fe
\end{parts}

\medskip

\noindent\textbf{Question 14.} AYUSH Kwatha is mainly useful in:\pts{1}
\begin{parts}
  \item Respiratory disorders and boosting immunity
  \item Eye diseases
  \item Bone disorders
  \item Skin burns
\end{parts}

\medskip

\noindent\textbf{Question 15.} Which is NOT a classical branch of Ashtanga Ayurveda?\pts{1}
\begin{parts}
  \item Radiology
  \item Kaya Chikitsa
  \item Shalya Tantra
  \item Kaumarbhritya
\end{parts}

\medskip

\noindent\textbf{Question 16.} The primary aim of Ayurveda is:\pts{1}
\begin{parts}
  \item Swasthasya Swasthya Rakshanam, Aturasya Vikara Prashamanam
  \item Wealth generation
  \item Cosmetic improvement
  \item None of the above
\end{parts}

\medskip

\noindent\textbf{Question 17.} Ahara (Diet), Nidra (Sleep), and Brahmacharya (Self-control) are called:\pts{1}
\begin{parts}
  \item Trayopastambha (Three Pillars of Life)
  \item Tridosha
  \item Trimala
  \item Triguna
\end{parts}

\medskip

\noindent\textbf{Question 18.} Dinacharya refers to:\pts{1}
\begin{parts}
  \item Daily health and hygiene routine
  \item Seasonal routine
  \item Night routine
  \item Dietary restrictions
\end{parts}

\medskip

\noindent\textbf{Question 19.} Ritucharya refers to:\pts{1}
\begin{parts}
  \item Seasonal health and diet regimen
  \item Daily routine
  \item Surgical rules
  \item Emergency care
\end{parts}

\medskip

\noindent\textbf{Question 20.} The concept of Agni in Ayurveda is primarily responsible for:\pts{1}
\begin{parts}
  \item Digestion, metabolism, and tissue transformation
  \item Circulation of blood
  \item Respiration
  \item Excretion of urine
\end{parts}

\medskip

\noindent\textbf{Question 21.} How many types of Agni are described in Ayurveda?\pts{1}
\begin{parts}
  \item 13 (1 Jatharagni, 5 Bhutagni, 7 Dhatvagni)
  \item 7
  \item 5
  \item 3
\end{parts}

\medskip

\noindent\textbf{Question 22.} Ojas is considered as the ultimate essence of:\pts{1}
\begin{parts}
  \item All seven Dhatus (Immunity and Vitality)
  \item Vata Dosha
  \item Blood waste
  \item Sweat
\end{parts}

\medskip

\noindent\textbf{Question 23.} Prakriti of an individual is established at the time of:\pts{1}
\begin{parts}
  \item Conception (Shukra-Shonita Samyoga)
  \item Birth
  \item Childhood
  \item Old age
\end{parts}

\medskip

\noindent\textbf{Question 24.} Panchakarma includes how many major purification therapies?\pts{1}
\begin{parts}
  \item 5 (Vamana, Virechana, Basti, Nasya, Raktamokshana)
  \item 3
  \item 7
  \item 8
\end{parts}

\medskip

\noindent\textbf{Question 25.} Abhyanga refers to:\pts{1}
\begin{parts}
  \item Therapeutic full-body oil massage
  \item Herbal steam
  \item Nasal drops
  \item Fasting
\end{parts}

\medskip

\noindent\textbf{Question 26.} Rasayana therapy aims at:\pts{1}
\begin{parts}
  \item Rejuvenation, longevity, and memory/immunity enhancement
  \item Immediate fever cure
  \item Surgical excision
  \item Bone setting
\end{parts}

\medskip

\noindent\textbf{Question 27.} Which Muhurta is considered best for waking up in the morning according to Ayurveda?\pts{1}
\begin{parts}
  \item Brahma Muhurta
  \item Nisha
  \item Rajani
  \item None of the above
\end{parts}

\medskip

\noindent\textbf{Question 28.} Which of the following is a Pranayama practice?\pts{1}
\begin{parts}
  \item Bhramari
  \item Bhujangasana
  \item Makarasana
  \item Shuttarungutasana
\end{parts}

\medskip

\noindent\textbf{Question 29.} Which of the following compendia is NOT in the group of Brihattrayi?\pts{1}
\begin{parts}
  \item Bhava Prakasha (belonging to Laghuttrayi)
  \item Charaka Samhita
  \item Sushruta Samhita
  \item Vagbhata Samhita
\end{parts}

\medskip

\noindent\textbf{Question 30.} ``Ayusho Vedah ....... is:''\pts{1}
\begin{parts}
  \item Ayurveda
  \item Yoga
  \item Samkhya
  \item Nyaya
\end{parts}

\vfill
\rule{\linewidth}{0.4pt}
\begin{center}\small
\href{https://bhu-exam.vercel.app/}{Click here to visit our website}\\[4pt]
\href{https://me-aryan.vercel.app/}{Click here to visit our contributor}
\end{center}

\end{document}
